\section{What is a quiver?}\label{sec:11-path-algebras:01-what-is-a-quiver:1}

A \textbf{quiver} is a \emph{directed graph}: a set of vertices and a set of directed edges (called \textbf{arrows}) between them. This is the same notion Mathlib calls \href{https://loogle.lean-lang.org/?q=Quiver}{\texttt{Quiver}} (built here from scratch, following Chapter 0's ``no Mathlib'' policy, rather than reusing Mathlib's). Formally, a quiver \(Q\) consists of:

\begin{itemize}
\tightlist
\item
  a set of vertices \(Q_0\),
\item
  a set of arrows \(Q_1\),
\item
  two functions \(s, t : Q_1 \to Q_0\) giving each arrow's \textbf{source} and \textbf{target}.
\end{itemize}

Picture an arrow \(\alpha : i \to j\) as a literal arrow drawn from vertex \(i\) to vertex \(j\); \(s(\alpha) = i\) and \(t(\alpha) = j\).

\subsection{Example quiver}\label{sec:11-path-algebras:01-what-is-a-quiver:2}

Take vertices \(\{1, 2, 3\}\) and arrows

\[
\alpha : 1 \to 2, \qquad \beta : 2 \to 3
\]
